\maketitle

\begin{abstract}
Long agent traces pack tool schemas, persistent identifiers, superseding
instructions, and filler into one key--value (KV) cache.  Every token has the
same storage cost, but not the same failure cost.  This motivates evicting by
\emph{application-visible structure}: protect schemas, constraints, sinks, and
the recent window; discard the rest.  We show that this works dramatically in one
regime, fails completely in another, and that a single measurable quantity
predicts which regime a workload is in.

On a synthetic stress benchmark (Qwen3-8B, matched 25\% keep), structure-aware
retention passes 112/120 examples against 1/120 for position-blind eviction, and
is statistically indistinguishable from reimplementations of SnapKV, PyramidKV,
and a structure--SnapKV hybrid at 108/120 (exact McNemar $p=0.125$).  On BFCL V3
multi-turn, scored by the unmodified official checker, the same policy scores
\textbf{0.000} on both Qwen3-8B ($n{=}140$) and Llama-3.1-8B ($n{=}143$), while
attention-based eviction remains indistinguishable from FullKV at the same
budget ($p=0.152$ and $p=1.0$).

The cause is measurable before any generation.  A role-based policy can rank only
while protected content is a \emph{minority} of context; above that it selects
everything and degenerates to index order.  Our synthetic benchmark is 94.9\%
filler (6.1\% protected); a BFCL system prompt \emph{is} 32 JSON tool schemas
(98.8\% protected).  We propose this protected fraction as a CPU-computable
pre-flight criterion, and ADAPT, a budget-relative prior that reduces to the
structure policy when it can rank and to attention selection when it cannot;
ADAPT ties SnapKV on BFCL ($p=1.0$) where the pure structure policy scores zero.

We map the operating envelope with the same rigor.  Structure matches FullKV
under mid-context relocation and degrades when state is buried in free-form
prose, bounding the prior to recognizable state.  Role-aware INT4 placement is
separately tested and shown not to be a quality lever, which keeps the storage
and retention claims independent.  The packed BF16/INT4 prototype reduces KV
payload to 0.719$\times$ at 1.200--1.215$\times$ time per output token, an
explicit space--time trade rather than an unqualified win.  We also publish a
correction to one released baseline (\S\ref{sec:random-defect}), so that every
column in this paper can be audited against the code.
\end{abstract}

\section{Introduction}
\label{sec:introduction}

Forty turns into an illustrative agent conversation, the model emits a tool call
without a required field.  The field belonged to a schema that the cache evicted
thirty turns earlier to stay under budget; byte for byte, it looked no different
from the filler that survived.  This is the failure PriorityKV targets.  In an
agent trace, a JSON schema field and a throwaway pleasantry cost the same storage
but not the same failure.  Losing the former can also resurrect an obsolete
instruction or substitute the wrong persistent identifier.  Figure~\ref{fig:failure-mode}
shows this running example schematically; it is illustrative, not an additional
measured trace.

Autoregressive inference stores prior keys and values so each decode step can
reuse the prefix.  The cache grows with sequence length and constrains serving
capacity~\cite{vaswani2017attention,kwon2023pagedattention}.  Existing work
reduces that cost through paging, streaming windows, attention-derived eviction,
or quantization~\cite{xiao2024streamingllm,zhang2023h2o,li2024snapkv,cai2024pyramidkv,liu2024kivi}.
Agent traces add a signal those methods do not directly use: their serialization
already exposes message roles, tool contracts, constraints, results, and state.

\begin{figure*}[!tbp]
  \centering
  \includegraphics[width=\linewidth]{agent_trace_failure_mode.pdf}
  \caption{\textbf{Equal storage cost, unequal failure cost.}  Two policies keep
  the same illustrative 14/26-token budget.  Hatched cells are evicted; callouts
  show the resulting failure modes rather than empirical counts.}
  \label{fig:failure-mode}
\end{figure*}

PriorityKV asks whether application-visible structure is useful when a cache
must discard tokens, and separately whether a mixed BF16/INT4 implementation
can save bytes.  These are different hypotheses: eviction removes information,
whereas quantization retains an approximation.  Combining their outcomes into
one ``compression quality'' claim would conceal the falsified INT4 hypothesis
and the measured latency cost.  Figure~\ref{fig:hypothesis-split} summarizes the
three experiment families and their distinct claim boundaries.

\begin{figure*}[!tbp]
  \centering
  \includegraphics[width=0.92\linewidth]{hypothesis_split.pdf}
  \caption{\textbf{Three questions, three verdicts.}  Eviction reliability,
  precision placement, and systems performance are evaluated independently;
  displayed counts and ratios are measured endpoints.}
  \label{fig:hypothesis-split}
\end{figure*}

\paragraph{Scope of claims.}
The paper separates three questions: whether structure helps eviction, whether
role-aware precision placement preserves quality, and whether packed storage
improves systems performance.  A result in one family does not establish either
of the others.  Our baselines are matched-budget repository reimplementations,
the benchmark is synthetic, and the systems measurements are single-request
prototype results.  We state additional limitations in
Section~\ref{sec:limitations} and do not repeat this boundary in every result.

This work makes four contributions.
\begin{enumerate}
  \item It defines PriorityBench-A, a deterministic synthetic benchmark for
  tool-schema conformance, instruction supersession, and multi-turn state recall,
  together with disjoint stress slices, placement controls, and a gold-span
  retention audit.
  \item It scales the decisive Qwen eviction comparison to $n=120$ and adds
  matched-budget attention baselines plus a paired McNemar test.  The supported
  claim is that structure matches or slightly exceeds SnapKV-class selection
  while decisively beating position-only baselines; it is not that structure
  beats FullKV or SnapKV in general.
  \item It reports two important negatives: structure loses to SnapKV on both
  Llama slices at a 5\% budget, and structure-aware placement of 75\% INT4
  positions does not create a quality advantage on the locked Qwen benchmark.
  \item It implements packed BF16/INT4 pages and two-call FlashInfer decode, then
  reports payload, allocator peak, packing, cold-scratch, end-to-end, and
  per-output-token measurements separately.
\end{enumerate}

\section{Related Work}
\label{sec:related}

\textbf{Paging and streaming.}
PagedAttention separates logical KV blocks from non-contiguous physical storage
to reduce fragmentation and enable sharing; it does not choose semantic tokens
to retain~\cite{kwon2023pagedattention}.  StreamingLLM identifies attention
sinks and combines initial tokens with a recent window for streaming language
modeling~\cite{xiao2024streamingllm}.  PriorityKV uses sink plus recent as a
mandatory region and as the shape of its position-only control, not as a claim
to reproduce StreamingLLM's streaming setting.

\textbf{Eviction and attention-derived selection.}
H2O retains accumulated-attention heavy hitters together with recent tokens
\cite{zhang2023h2o}.  SnapKV uses an observation window near the end of the
prompt to select clustered prefix features~\cite{li2024snapkv}.  PyramidKV varies
the cache budget across layers in response to pyramidal information funneling
\cite{cai2024pyramidkv}.  These methods infer importance from model behavior;
PriorityKV tests application-visible message structure.  The attention-baseline
comparison uses repository reimplementations at matched token budgets, not the
authors' reference code.  In particular, H2O is a chunked SDPA
reimplementation and carries an explicit fidelity caveat.  A hybrid that
protects structural tokens and fills the remainder with SnapKV was evaluated:
it equals SnapKV on Qwen at 25\% keep and collapses below both parents on the two
Llama 5\% slices.  Thus the proposed complementarity hypothesis is not
supported.

\textbf{KV quantization and attention engines.}
KIVI motivates asymmetric group-wise low-bit KV quantization and keeps residual
tokens at full precision~\cite{liu2024kivi}.  PriorityKV's INT4 path is not a new
quantizer and is not positioned as a replacement for KIVI.  It is a real packed
storage path for testing role-aware precision placement and for accounting
payload bytes.  FlashInfer exposes composable attention kernels and
log-sum-exp state merging for heterogeneous KV layouts~\cite{ye2025flashinfer}.
PriorityKV uses that merge contract, while retaining an explicit BF16 cold
scratch limitation.

\textbf{Long-context evaluation.}
LongBench and RULER test broad long-context capabilities
\cite{bai2024longbench,hsieh2024ruler}.  Our suite instead targets exact failure
attribution in synthetic agent protocols.  It cannot establish workload
prevalence or replace general long-context suites.

\section{PriorityBench-A}
\label{sec:benchmark}

\subsection{Tasks and scoring}
\label{sec:tasks}

PriorityBench-A generates multi-turn chats in which an early span contains the
state required by a final request and intervening turns provide benign filler.
Scoring is binary and programmatic; no LLM judge is used.  Table~\ref{tab:tasks}
summarizes the three equally represented categories in the locked 240-example
manifest.

\begin{table*}[!tbp]
  \centering
  \caption{PriorityBench-A task families in the locked manifest ($n=240$;
  80 examples per category).}
  \label{tab:tasks}
  \begin{tabular}{lll}
    \toprule
    Category & Required behavior & Deterministic scorer \\
    \midrule
    Tool schema & Emit a valid call under an early contract & JSON/schema subset \\
    Instruction supersession & Follow the latest conflicting constraint & Required/forbidden regex \\
    Multi-turn state & Reuse an earlier ID, path, or preference & Exact slot match \\
    \bottomrule
  \end{tabular}
\end{table*}

Two manifests from the same PriorityBench-A generator serve different studies.
The locked 240-example manifest (W3) contains 83, 81, and 76 examples at nominal
8k, 16k, and 32k context strata; stable hashing assigns 92 calibration, 49
validation, and 99 test examples.  The 120-example stress manifest (W5) is split
evenly across the three categories and 8k/16k strata, with three disjoint
40-example slices.  Reproducibility metadata, including both manifest digests,
appears in Table~\ref{tab:artifacts} rather than in the narrative.

\subsection{Placement and leakage controls}
\label{sec:controls}

The generator's visible templates are also a threat: a heuristic tagger can
appear semantic if gold state is consistently short, early, or marked by a
template recognizable to the policy.  We therefore evaluate two transformations
of each stress slice.  \emph{Middle relocation} moves state away from the prefix
while retaining the leading system message and final request.  \emph{Buried
state} lengthens or embeds state-bearing turns so that a short-span heuristic
does not protect all gold.

A separate CPU audit maps the gold-bearing messages into tokenizer spans and
intersects them with the selected positions.  This audit uses no model outputs.
It measures the gold fraction already in sink plus recent and the gold fraction
retained by each policy.  Thus placement robustness and label leakage into the
always-kept window are tested independently.

\section{Method}
\label{sec:method}

\subsection{Roles and matched-budget retention}
\label{sec:retention-method}

Let a tokenized trace be $x_{1:n}$ with structural roles $r_{1:n}$.  For keep
fraction $k$, the nominal budget is
\begin{equation}
  B(n,k)=\operatorname{round}(kn).
  \label{eq:budget}
\end{equation}
The mandatory set is $M=M_{\mathrm{sink}}\cup M_{\mathrm{recent}}$, where the
first 16 tokens and last 128 tokens are protected.  In the evaluated token-level
regime, each compressed arm selects $A\supseteq M$ with
\begin{equation}
  |A|=\max(B(n,k),|M|).
  \label{eq:matched}
\end{equation}
Uniform retention uses the mandatory sink and fills from the recent edge.
Random retention keeps $M$ and samples the remaining budget.  Structure-aware
retention first includes system, tool, constraint, and other state-bearing
positions, then fills any residual budget from the recent edge.  Whole-page
experiments apply the same priority order while respecting page boundaries and
never exceeding the token budget except for $M$.

The tagger consumes message roles and transparent schema or constraint markers.
It does not inspect benchmark answers or future attention.  The policy protects
explicit protocol roles more reliably than unmarked free-form state; the buried
controls in Section~\ref{sec:p0-results} quantify that boundary.

\subsection{Pages and mixed precision}
\label{sec:mixed-method}

The physical page size is $P=16$ tokens.  In the mixed-precision study, every
position remains present and an exact target fraction $f=0.75$ is stored as
INT4, subject to the common sink/recent BF16 protection.  The role-blind arm
spaces INT4 positions across the demotable region.  The structure arm demotes
filler, generated, and low-risk positions first, spilling into protected soft
roles only if the byte budget requires it.  Both arms therefore have the same
number of INT4 positions.

Figure~\ref{fig:allocation} traces the implemented path from tagging to the page
table.  It separates policy decisions---retention for eviction or dtype for
mixed precision---from physical payload.  Cold K and V pages contain packed
4-bit values plus per-group FP32 scale metadata; hot pages remain BF16.

\begin{figure*}[!tbp]
  \centering
  \includegraphics[width=\linewidth]{page_allocation_architecture.pdf}
  \caption{\textbf{PriorityKV architecture.}  Chat metadata becomes token roles
  and 16-token pages before matched-budget retention or precision placement.
  The page records are illustrative.}
  \label{fig:allocation}
\end{figure*}

\subsection{Two-call FlashInfer decode}
\label{sec:decode-method}

Prefill uses the native Hugging Face path.  The resulting cache is partitioned
into coalesced hot BF16 pages and packed cold pages.  Before the frozen full-scratch
decode measurement, cold pages are dequantized into BF16 GPU scratch.  Each
layer makes at most two homogeneous FlashInfer calls: one over hot pages plus
the decode tail and one over cold scratch.  Their outputs and native
log-sum-exp states are combined with \texttt{merge\_state}.  A parity artifact
records maximum absolute error of approximately $4.88\times10^{-4}$ after
correcting an LSE-base mismatch.

Figure~\ref{fig:lifetime} shows why cache payload, allocator peak, and latency
must be reported separately.  Packed bytes persist, full cold scratch is live
during decode, and the decode tail grows token by token.  The streamed-cold
experiment (P2) instead dequantizes chunks and releases them after merging, but its
three-example run is only a successful smoke test and is not a comparative
systems result.

\begin{figure*}[!tbp]
  \centering
  \includegraphics[width=\linewidth]{decode_memory_lifetime.pdf}
  \caption{\textbf{Decode path and allocation lifetime.}  Two homogeneous
  FlashInfer calls are merged through log-sum-exp state; packed payload persists
  while BF16 cold scratch is live during decode.}
  \label{fig:lifetime}
\end{figure*}

\section{Experimental Design}
\label{sec:design}

All GPU experiments use greedy decoding on NVIDIA H200 hardware.  Canonical
jobs use at most two GPUs, while each eviction, attention-baseline, and transfer
job uses one.  Table~\ref{tab:artifacts} records the pinned model revisions,
manifest digests, and result bundles.

\subsection{Eviction, attention baselines, and placement controls}
\label{sec:p0-design}
\label{sec:p1-design}

The eviction and placement study (P0) evaluates Qwen on the 120-example stress
manifest, divided into three disjoint 40-example slices and balanced across 8k
and 16k contexts.  At token keep $k=0.25$, it compares structure, uniform,
random, keep-all, and a separately recorded FullKV run using deterministic task
pass rate.  Middle-relocation and buried-state transformations are evaluated on
all three slices.  The attention-baseline comparison (P1) reuses the same
examples, model, contexts, and budget for SnapKV, PyramidKV, H2O, and a
structure--SnapKV hybrid.  SnapKV uses a 64-token observation window and kernel
size five.  The H2O arm is a chunked SDPA reimplementation and carries a fidelity
limitation.  A tracked artifact supplies the exact two-sided paired McNemar test
for structure versus SnapKV.

\subsection{CPU retention audit}
\label{sec:audit-design}

The audit runs the Qwen and Llama tokenizers on the three stress slices for each
model ($n=120$ each), maps gold messages to token spans, and applies uniform,
random, and structure selection at $k=0.25$.  It reports mean token-retention
fractions rather than model accuracy.  These are descriptive span means, so no
binomial interval or significance test is attached.

\subsection{Mixed-precision quality and systems measurements}
\label{sec:int4-design}
\label{sec:systems-design}

The locked mixed-precision study evaluates Qwen on all 240 examples: 83 at 8k,
81 at 16k, and 76 at 32k.  FullKV, uniform-mixed, and structure-mixed arms use
matched $f=0.75$ INT4 placement for the mixed arms.  The H200 packed FlashInfer
path is scored by the same deterministic benchmark metric.  This family tests a
precision-placement quality hypothesis, not eviction.  No paired significance
artifact exists; we report Wilson intervals from the observed success counts.

The latency measurement uses nine 8k and nine 16k examples, 128 generated
tokens, one warm-up, and three timing repetitions summarized by per-example
medians.  FullKV uses SDPA while the packed path uses the FlashInfer shim; this
is an end-to-end prototype comparison, not a kernel-isolated study.  The memory
measurement uses 18 examples and records packed payload plus allocated and
reserved CUDA peak.  Both are single-request studies without concurrent serving
or tail latency.  Canonical jobs were not independently repeated, so no error
bars are shown.  A three-example streamed-cold run (P2) checks exit status and
parity but serves only as a smoke test.

\subsection{Cross-model transfer and weak checks}
\label{sec:p3-design}
\label{sec:legacy-design}

The Llama transfer study (P3) evaluates structure, SnapKV, H2O, PyramidKV, and
hybrid on all three stress slices at $k=0.25$ ($n=120$).  At $k=0.05$, only the
first two slices are run ($n=40$ each); no uniform Llama arm is available.  Two
earlier checks remain deliberately secondary: one fixed 14-example Qwen slice
is reused across three whole-page budgets, and a 14-example Gemma-2-9B-IT check
uses Qwen-tuned prompts and has a low FullKV score.

\section{Results}
\label{sec:results}

\subsection{Structure protects state that blind eviction drops}
\label{sec:p0-results}

At matched 25\% token keep, structure passes 112/120 Qwen examples while the
position-blind control passes 1/120 (Table~\ref{tab:qwen-main}).  The
corresponding Wilson intervals do not come close to overlapping.  We report a
single position-blind control rather than two: our released ``random'' arm is
byte-identical to the uniform arm, a defect characterized in
\S\ref{sec:random-defect} rather than left for a reader to discover.  Figure~\ref{fig:eviction}
places the position-blind and attention-derived selectors on one pass-rate axis.

\begin{figure*}[!tbp]
  \centering
  \includegraphics[width=0.88\linewidth]{eviction_and_baselines.pdf}
  \caption{\textbf{Matched-budget Qwen comparison} ($n=120$, $k=0.25$).
  Points show pass rates with Wilson 95\% intervals; exact counts are aligned at
  right and the bracket reports the paired structure--SnapKV McNemar test.
  H2O* denotes the chunked repository reimplementation.}
  \label{fig:eviction}
\end{figure*}

\begin{table}[!htbp]
  \centering
  \caption{Pooled Qwen eviction and attention-baseline results at $k=0.25$
  ($n=120$).  Intervals are Wilson 95\%.}
  \label{tab:qwen-main}
  \begin{tabular}{lrrl}
    \toprule
    Arm & Passes & Rate & Wilson 95\% interval \\
    \midrule
    FullKV & 107/120 & 0.892 & [0.823, 0.936] \\
    Structure & 112/120 & 0.933 & [0.874, 0.966] \\
    Position-blind$^{\dagger}$ & 1/120 & 0.008 & [0.002, 0.046] \\
    \midrule
    SnapKV & 108/120 & 0.900 & [0.833, 0.942] \\
    PyramidKV & 108/120 & 0.900 & [0.833, 0.942] \\
    Hybrid & 108/120 & 0.900 & [0.833, 0.942] \\
    H2O (chunked) & 82/120 & 0.683 & [0.596, 0.760] \\
    \bottomrule
  \end{tabular}

  \vspace{0.35em}
  \parbox{0.86\linewidth}{\footnotesize $^{\dagger}$Earlier versions of this
  table reported ``uniform'' and ``random'' as two arms, both at 1/120.  They are
  a single computation: our \texttt{select\_random} is byte-identical to
  \texttt{select\_uniform} at every context length (\S\ref{sec:random-defect}).
  We report one row rather than double-count one control.}
\end{table}

The CPU audit confirms that this gap reflects retention rather than an easy
always-kept window.  Only 1.0--1.3\% of Qwen gold tokens and none of the Llama
gold tokens lie in sink plus recent (Table~\ref{tab:audit}).  Structure retains
all mapped gold tokens, whereas uniform retains only the small fraction already
in that window.

\begin{table}[!htbp]
  \centering
  \caption{Gold-span retention audit at $k=0.25$ ($n=40$ per slice).  Entries
  are mean token fractions, not pass counts.}
  \label{tab:audit}
  \begin{tabular}{llrrr}
    \toprule
    Model & Slice & Gold in sink+recent & Structure kept & Uniform kept \\
    \midrule
    Qwen & 0 & 0.010 & 1.000 & 0.010 \\
    Qwen & 1 & 0.013 & 1.000 & 0.013 \\
    Qwen & 2 & 0.010 & 1.000 & 0.010 \\
    Llama & 0 & 0.000 & 1.000 & 0.000 \\
    Llama & 1 & 0.000 & 1.000 & 0.000 \\
    Llama & 2 & 0.000 & 1.000 & 0.000 \\
    \bottomrule
  \end{tabular}
\end{table}

Placement controls bound the positive result (Table~\ref{tab:placement}).
Moving state to the middle gives both structure and FullKV 115/120, or 0.958
[0.906, 0.982].  Burying state in longer free-form turns lowers structure to
80/120, or 0.667 [0.578, 0.745], versus FullKV at 107/120, or 0.892
[0.823, 0.936].  Uniform scores 3/120 in the middle condition and 0/120 when
state is buried.  Visible roles are therefore a useful prior, not an oracle for
arbitrarily embedded state.

\begin{table}[!htbp]
  \centering
  \caption{Placement controls ($n=40$ per slice), reported as structure / FullKV
  / uniform pass rates.}
  \label{tab:placement}
  \begin{tabular}{lcc}
    \toprule
    Slice & Middle relocation & Buried state \\
    \midrule
    0 & 0.975 / 0.975 / 0.025 & 0.675 / 0.900 / 0.000 \\
    1 & 0.950 / 0.950 / 0.050 & 0.650 / 0.875 / 0.000 \\
    2 & 0.950 / 0.950 / 0.000 & 0.675 / 0.900 / 0.000 \\
    \bottomrule
  \end{tabular}
\end{table}

\subsection{Structure matches attention selection on this benchmark; the hybrid does not help}
\label{sec:p1-results}

\emph{Scope note.}  The parity reported here is specific to PriorityBench-A and
does not survive transfer.  \S\ref{sec:external} measures the same two policies
on an externally authored benchmark and finds attention selection intact while
the structure policy collapses to zero.  The two results are consistent once the
protected fraction is measured (\S\ref{sec:external-boundary}); readers should
carry the qualifier, not the headline.

Structure passes 112/120, while SnapKV, PyramidKV, and hybrid each pass 108/120
(Table~\ref{tab:qwen-main}).  The structure interval overlaps every 0.900 arm.
The paired structure--SnapKV artifact records $b=4$ cases where
structure alone passes and $c=0$ in the opposite direction; exact two-sided
McNemar $p=0.125$.  The difference is not statistically significant.  Hybrid
equals SnapKV rather than exceeding it, falsifying the complementarity
hypothesis.  H2O passes 82/120 using the chunked implementation;
its implementation caveat prevents a reference-H2O claim.

\subsection{The advantage depends on model and budget}
\label{sec:audit-results}
\label{sec:p3-results}

All five tested structure and attention arms pass all 120 Llama examples at
$k=0.25$; the Wilson interval for 120/120 is [0.969, 1.000].  With no uniform
arm, this is an easy-task ceiling rather than evidence of universal transfer.
At $k=0.05$, SnapKV passes 40/40 on each evaluated slice, while structure passes
35/40 and 36/40.  No paired test was produced for these Llama slices.

\begin{figure*}[!tbp]
  \centering
  \includegraphics[width=0.92\linewidth]{budget_and_transfer.pdf}
  \caption{\textbf{Budget and model dependence.}  Qwen and Llama at $k=0.25$
  use $n=120$; each Llama $k=0.05$ condition uses $n=40$.  Error bars are Wilson
  95\% intervals.}
  \label{fig:transfer}
\end{figure*}

At the tighter budget, failures concentrate in tool-schema tasks.  Structure
passes 8/13 and 9/13 schema examples in the first and second slices, while
SnapKV passes 13/13 in both.  Hybrid is worse at 1/13 and 0/13, with overall
rates 0.575 and 0.525.  Protecting broad structural spans can therefore consume
too much residual budget, and simply unioning role and attention selections does
not make them complementary.

\subsection{Matched INT4 placement does not separate quality}
\label{sec:int4-results}

On the locked 240 examples, FullKV passes 213, uniform-mixed 211, and
structure-mixed 212 (Table~\ref{tab:int4}).  The structure--uniform difference is
one example and no paired statistical artifact supports a quality distinction.
At 8k and 16k, every arm is 1.0; at 32k, FullKV, uniform, and structure are
49/76, 47/76, and 48/76.  Figure~\ref{fig:int4} shows the count-backed Wilson
intervals by context.  The tested role-aware INT4 quality hypothesis is
falsified.

\begin{figure*}[!tbp]
  \centering
  \includegraphics[width=0.88\linewidth]{lock240_quality_by_length.pdf}
  \caption{\textbf{Matched INT4 quality.}  Qwen3-8B with FullKV or matched
  $f=0.75$ INT4 placement.  Error bars are Wilson 95\% intervals from exact
  counts (83, 81, and 76 examples).}
  \label{fig:int4}
\end{figure*}

\begin{table}[!htbp]
  \centering
  \caption{Locked Qwen quality ($n=240$).  Intervals are Wilson 95\%.}
  \label{tab:int4}
  \begin{tabular}{lrrl}
    \toprule
    Arm & Passes & Rate & Wilson 95\% interval \\
    \midrule
    FullKV & 213/240 & 0.8875 & [0.841, 0.922] \\
    Uniform mixed & 211/240 & 0.8792 & [0.832, 0.915] \\
    Structure mixed & 212/240 & 0.8833 & [0.837, 0.918] \\
    \bottomrule
  \end{tabular}
\end{table}

\subsection{Packed bytes incur scratch and latency costs}
\label{sec:systems-results}

The mixed structure path has a modeled ratio of 0.473$\times$, measured packed
payload ratio of 0.719$\times$, and allocated-peak ratio of 0.868$\times$ FullKV
on the 18-example memory slice.  The gap is caused by BF16 cold scratch, not by
an accounting inconsistency.  Figure~\ref{fig:systems} reports memory and latency
ratios on the same parity reference.  No uncertainty bars are shown because the
canonical jobs were not independently repeated.

\begin{figure*}[!tbp]
  \centering
  \includegraphics[width=0.88\linewidth]{systems_tradeoff.pdf}
  \caption{\textbf{H200 systems trade-off.}  Packed payload is smaller than
  FullKV, allocator peak falls less, and every measured latency ratio exceeds
  the dashed 1.0 reference.}
  \label{fig:systems}
\end{figure*}

At 8k, structure packing and scratch construction take 34.8 and 14.2 ms.  E2E
TTFT is 1.118$\times$ FullKV, and TPOT rises from 25.25 to 30.29 ms
(1.200$\times$).  At 16k, pack and scratch take 48.1 and 19.7 ms; E2E is
1.113$\times$, and TPOT rises from 23.84 to 28.95 ms (1.215$\times$).  These are
regressions, not speedups.  The streamed-cold smoke run exits successfully with parity and
approximately 36.4 GiB peak in its log, but with only three examples and no
comparative frontier it remains a smoke test rather than a systems result.

\subsection{Additional weak checks}
\label{sec:legacy-results}

A whole-page Qwen ablation reuses the same 14 examples at 15\%, 25\%, and 35\%
keep.  Structure remains 9/14 at every budget; random rises from 1/14 to 6/14,
and the role-blind arm remains 0/14.  Because the slice is reused, this is one
small ablation rather than three replications.  On a separate 14-example Gemma
check, FullKV passes 5/14, structure 2/14, and uniform 0/14.  The low FullKV
score, Qwen-tuned prompts, and sample size make both results directional checks,
not transfer evidence.

\section{External Evaluation}
\label{sec:external}

The results above are measured on PriorityBench-A, a benchmark we constructed.
This section tests whether the eviction prior transfers to a benchmark we did
not author.  We evaluate on BFCL V3 multi-turn using the unmodified official
\texttt{multi\_turn\_checker}, which executes tool calls against stateful API
implementations and compares final state and execution responses against ground
truth.  All non-FullKV arms are \texttt{kvpress} presses over an identical full
prefill at an identical compression ratio, so the arms differ only in which KV
entries survive, never in how they are removed.

\subsection{The structural signal is exhausted in schema-dominated contexts}
\label{sec:external-transfer}

Table~\ref{tab:external-arms} reports 150 conversations per model at a 25\% keep
budget, scored per conversation.

\begin{table}[!htbp]
  \centering
  \caption{\textbf{BFCL V3 multi-turn, official scorer, 25\% keep.}  Paired
  completeness 0.933 (Qwen) and 0.953 (Llama); exclusions are all
  \texttt{MODEL\_CONTEXT\_LIMIT}; zero matched-budget violations.}
  \label{tab:external-arms}
  \begin{tabular}{lrr}
    \toprule
    Arm & Qwen3-8B ($n{=}140$) & Llama-3.1-8B ($n{=}143$) \\
    \midrule
    FullKV            & 0.193 & 0.077 \\
    SnapKV            & 0.136 & 0.084 \\
    ADAPT (\S\ref{sec:external-adapt}) & 0.129 & --- \\
    Structure         & 0.000 & 0.000 \\
    Uniform           & 0.000 & 0.000 \\
    Random            & 0.000 & 0.000 \\
    \bottomrule
  \end{tabular}
\end{table}

\textbf{Is the harness sound?}  FullKV itself reaches only 0.193 and 0.077, so a
reader should ask whether the pipeline is broken rather than the policy.  Four
checks argue otherwise.  (i)~Scoring is the unmodified upstream checker, and
replaying the dataset's own ground truth through our rollout scores \emph{valid},
while corrupted arguments and dropped turns score invalid.  (ii)~BFCL V3
multi-turn is known-hard for 8B models in prompting mode; low absolute accuracy
is expected at this scale.  (iii)~The arms separate far above the floor: a Wilson
upper bound of 0.027 on the zero arms excludes the FullKV rate on both models.
(iv)~An earlier configuration \emph{did} return 0.000 for every arm including
FullKV; we traced it to reasoning blocks reaching the decoder unstripped, fixed
it, and FullKV rose to a non-degenerate rate.  We report the diagnostic because
it is the exact failure mode a skeptical reader should suspect.

Structure-aware retention scores zero on every conversation, on both models.
Exact paired McNemar gives FullKV versus structure $p=1.5\times10^{-8}$ (Qwen)
and $p=9.8\times10^{-4}$ (Llama).  Attention selection, by contrast, is
statistically indistinguishable from FullKV at a $4\times$ budget: $p=0.152$ on
Qwen and $p=1.0$ on Llama.  Structure versus SnapKV is decisive against
structure on both ($p=3.8\times10^{-6}$ and $p=4.9\times10^{-4}$).  Two
architectures reproduce every sign.

\subsection{The protected fraction predicts the outcome}
\label{sec:external-boundary}

The failure has a measurable cause.  A role-based policy can only \emph{rank}
while the protected mass is smaller than the keep budget; above that it selects
everything and degenerates to index order.  Figure~\ref{fig:protected-fraction}
measures that quantity per workload with the same tagger used throughout.

\begin{figure*}[!tbp]
  \centering
  \includegraphics[width=\linewidth]{protected_fraction_boundary.pdf}
  \caption{\textbf{Protected mass relative to the keep budget.}  PriorityBench-A
  is 94.9\% filler, so 6.1\% of tokens carry a protected role and the policy has
  ample room to rank within a 25\% budget.  A BFCL system prompt \emph{is} 32
  JSON tool schemas, so 98.8\% of tokens carry the protected \textsc{tool} role
  and nothing is available to discard.  $\tau$-bench sits between at 81.6\%; it
  is a retention-only audit and has no task-success value.}
  \label{fig:protected-fraction}
\end{figure*}

\begin{figure*}[!tbp]
  \centering
  \includegraphics[width=\linewidth]{external_bfcl_arms.pdf}
  \caption{\textbf{Paired BFCL accuracy with Wilson intervals.}  The same
  ordering holds on both models: FullKV and attention selection are
  indistinguishable, while the three non-attention policies are at zero with a
  Wilson upper bound of 0.027.}
  \label{fig:external-arms}
\end{figure*}

This bounds the PriorityBench-A result rather than contradicting it.  The
synthetic benchmark contains distractor filler by construction, which is exactly
the regime in which ``protect the structure, drop the filler'' is a winning rule.
Real tool-calling prompts contain no such filler.

A generation-free audit of 4{,}856 public $\tau$-bench trajectories
(828{,}000 extracted spans) corroborates the mechanism from the other side:
structure retains explicit policy lines at 0.820 versus 0.001 for a position-only
budget, but loses on reused identifiers (0.055 versus 0.315).  Structure wins
precisely where protected spans are scarce and separable.

\subsection{ADAPT: structure as a budget-relative prior}
\label{sec:external-adapt}

The hybrid falsified in \S\ref{sec:legacy-results} force-protects structural
positions, which consumes the entire budget whenever protected mass exceeds it.
That is a mis-specification rather than a dead end.  ADAPT replaces the hard
constraint with a budget-relative prior,
\begin{align}
  \alpha &= \min\!\left(1,\ \frac{\text{keep budget}}{\text{protected mass}}\right),
  \label{eq:adapt-alpha}\\[2pt]
  s &= \alpha\,\mathrm{rank}(s_{\text{struct}})
     + (1-\alpha)\,\mathrm{rank}(s_{\text{attn}}).
  \label{eq:adapt}
\end{align}
where both quantities are known from the prompt, so $\alpha$ involves no fitting
and no free parameter.  Rank normalization is required because the structure and
attention scores differ by several orders of magnitude.  At $\alpha=1$ the press
provably selects exactly what the structure arm selects, so ADAPT subsumes the
policy of \S\ref{sec:method} rather than approximating it.

On BFCL the measured $\alpha$ is $0.267$ on average (range $0.250$--$0.401$) over
833 generation steps, against $\approx 0.25$ implied by the 98.8\% protected
fraction.  ADAPT ties SnapKV (exact McNemar $p=1.0$, $\Delta=-0.007$, 95\% CI
$[-0.057, +0.071]$) and is indistinguishable from FullKV ($p=0.108$), while
beating structure, uniform, and random ($p=7.6\times10^{-6}$).  We therefore
claim only that ADAPT is never worse and recovers attention-level behavior
automatically from a measurement; we do \emph{not} claim it beats SnapKV.  The
$\alpha$ expression was written before the ADAPT run but committed afterwards, so
the prediction is pre-specified by construction rather than by timestamp.

\section{Discussion}
\label{sec:discussion}

\textbf{Structure is a prior, not an oracle.}
The eviction results and retention audit jointly show that visible protocol structure can
protect state that a sink/recent or random budget discards.  The middle and
buried controls give the necessary qualification: when relevant state remains
recognizable, structure ties FullKV; when it is buried inside longer free-form
turns, structure loses to FullKV.  A production policy would need richer state
metadata or a learned signal for unmarked spans.

\textbf{The advantage depends on model and budget.}
Qwen at 25\% keep strongly separates structure from position-only policies and
places structure within four paired successes of SnapKV.  Llama at 25\% cannot
differentiate any competent selector; at 5\%, SnapKV is better on both slices.
The tool-schema drop at that extreme budget suggests that role protection alone
can allocate too much or too little within a broad protected span, while
attention can choose a smaller subset.  This is a useful crossover, not an
anomaly to hide.

\textbf{Hybrid complementarity is falsified at the tested settings.}
On Qwen, hybrid equals SnapKV; on Llama at 5\%, it is below both parents.  A union
of protected positions with attention-ranked positions changes the residual
budget and can inherit the weaknesses of both rules.  Any future hybrid claim
requires a new allocation rule and new evidence rather than reinterpretation of
the current comparison.  \S\ref{sec:external-adapt} supplies exactly that: the
failure is that a hard union consumes the whole budget once protected mass
exceeds it, and replacing the union with a budget-relative weight removes the
pathology.  We present ADAPT as a distinct rule tested on new data, not as a
reinterpretation of the falsified hybrid.

\textbf{The protected fraction is the operative variable.}
The external evaluation (\S\ref{sec:external}) turns the model-and-budget
dependence above into a single measurable quantity.  Structure-aware retention
helps when protected content is a minority of context and fails when it is the
majority, because a binary role label cannot rank within a protected set that
already exceeds the budget.  The fraction is computable on CPU before any
generation, which makes it a usable pre-flight check rather than a post-hoc
explanation.

\textbf{Bytes, peak, and speed are different systems quantities.}
Packed values and scale metadata genuinely reduce payload.  Full-layer BF16
cold scratch then expands the live working set and adds preparation and decode
cost.  The observed 0.719$\times$ payload, 0.868$\times$ peak, and roughly
1.20$\times$ TPOT are mutually consistent.  They do not establish serving
throughput or capacity under concurrency.

\section{Limitations and Threats to Validity}
\label{sec:limitations}

\textbf{Synthetic workload.}
PriorityBench-A is generated from controlled templates and exact scorers.  It
measures whether particular protocol state survives, not the prevalence of
these failures in deployed agent traffic, and it is not a LongBench or RULER
matrix.  \S\ref{sec:external} converts this from a caveat into a measurement:
the benchmark's 94.9\% filler places it at one end of the protected-fraction
axis, and the structure advantage does not survive at the other end.

\textbf{External evaluation scope.}
The external results cover one budget (25\% keep), two 8B models, and 150
conversations per model.  BFCL V3 multi-turn is hard for models of this size, so
FullKV itself scores 0.193 and 0.077; the arms are separated well above that
floor, but absolute headroom is small.  The $\tau$-bench audit measures retention
only---no user simulator, no tool backend, no generation---and excludes SnapKV,
whose scores require realized attention.

\textbf{Tagger--generator alignment.}
The tagger uses role, length, and schema/constraint markers that resemble the
generator's visible templates.  The gold audit rules out answer labels already
in sink plus recent, but it does not eliminate all template alignment.  The
buried result---0.675, 0.650, and 0.675 versus FullKV 0.900, 0.875, and
0.900---is the clearest measured boundary.

\subsection{Correction to a released position-blind baseline}
\label{sec:random-defect}

Auditing our own released code against the paper surfaced one reporting error,
which we correct here.  The synthetic study listed ``uniform'' and ``random'' as
two position-blind controls; they are one computation.  In \texttt{select\_random} the forced
window is sized as $\text{recent} = \max(\text{force\_recent},\,
\text{budget} - \text{sink})$; whenever the budget exceeds $\text{sink} +
\text{force\_recent}$ this makes the forced sink-plus-recent block exactly fill
the budget, the residual is zero, and the sampling branch never executes.  We
verified byte-identical index sets at $n \in \{600,\,4096,\,40000\}$, with zero
tokens ever drawn at random.

The consequence is contained.  No measured value changes, because the column was
always reporting the uniform policy correctly; what changes is the count of
independent controls behind the position-blind comparison, from two to one.  The
structure-versus-position-blind separation---112/120 against 1/120---is unaffected
in magnitude and significance.  Our external evaluation uses a corrected
position-blind arm that holds sink and recent at their configured sizes and
samples the remainder; it also scores 0.000 on BFCL, so \S\ref{sec:external}
stands as reported.  We publish the correction because reproducibility claims are
only worth the audits behind them.

\textbf{Baseline fidelity.}
SnapKV, PyramidKV, H2O, and hybrid are matched-budget repository
reimplementations, not runs of reference implementations.  Chunked H2O under
SDPA is particularly approximate.  The paper therefore does not claim a
reference-code leaderboard comparison.

\textbf{Statistical scope.}
The paired Qwen structure--SnapKV result is not significant ($p=0.125$).  The
eviction, baseline, and transfer studies are single canonical runs over
deterministic example sets, not independent
machine replications.  Wilson intervals describe finite pass counts; they do
not quantify hardware-run variability.  The Llama 5\% comparison has no tracked
paired test.

\textbf{Model and budget scope.}
Qwen3-8B on H200 is the primary matrix.  Llama provides a ceiling and a negative
crossover, not universal transfer.  Gemma is a reduced 14-example legacy check
with a low FullKV score.  Only selected 5\%, 25\%, and legacy page budgets are
tested.

\textbf{Systems scope.}
The latency comparison uses different backends (FullKV SDPA and mixed
FlashInfer) and single-request measurements.  It does not measure batching,
concurrent serving, throughput, or tails.  Payload is not peak memory because
cold attention expands INT4 into BF16.  The streamed-cold variant records exit zero,
parity, and a log peak, but it has no FullKV/TPOT frontier and is not a systems
result.

\section{Reproducibility}
\label{sec:reproducibility}

The repository freezes model revisions, manifests, configuration files, result
bundles, and job commands.  CPU logic and figures can be regenerated without an
H200; GPU re-execution requires the pinned model weights and isolated GPU
environment described in \texttt{docs/REPRODUCIBILITY.md}.  Table~\ref{tab:artifacts}
maps every headline claim to its tracked artifact.

\begingroup
\footnotesize
\sloppy
\hbadness=10000
\begin{table*}[!tbp]
  \centering
  \caption{Reproducibility index.  Result paths are relative to
  \texttt{jobs/results/}; exact job commands appear in the final run manifest.}
  \label{tab:artifacts}
  \begin{tabular}{@{}p{0.27\linewidth}p{0.66\linewidth}@{}}
    \toprule
    Item or claim & Pinned revision, digest, or result artifact \\
    \midrule
    Qwen3-8B revision & \texttt{b968826d9c46dd6066d109eabc6255188de91218} \\
    Llama-3.1-8B revision & \texttt{0e9e39f249a16976918f6564b8830bc894c89659} \\
    Locked manifest (240) & SHA256 \texttt{fc44b966725738c94008ba61ce57ad73}\allowbreak\texttt{66169b9c0be73074f8161d909ccfae89} \\
    Stress manifest (120) & SHA256 \texttt{154c07b900a7be6a26a0381b0fa28d49}\allowbreak\texttt{e0847086eade0b98e9472663b578784f} \\
    Qwen eviction & Three \texttt{p0\_w5\_s*\_kf25\_token\_*/summary.json} files \\
    Placement controls & Six \texttt{p0[a--f]\_w5\_*\_\{middle,buried\}\_*/summary.json} files \\
    Attention baselines & Three \texttt{p1\_attn\_baselines\_s*\_kf25\_*/summary.json} files \\
    H2O 82/120 & \texttt{p1\_h2o\_chunked\_s0\_kf25\_gpu1\_r1/summary.json} plus baseline summaries for slices 1--2 \\
    McNemar $p=0.125$ & \texttt{p1\_structure\_vs\_snapkv\_mcnemar.json} \\
    Gold-span audit & Six \texttt{audit\_retention\_*\_kf25\_summary.json} files \\
    Llama transfer & \texttt{p3\_llama31\_attn\_s*\_kf\{25,05\}\_*/summary.json} \\
    Locked INT4 quality & \texttt{mg\_b\_lock240\_quality\_gpu01\_r1/summary.json} \\
    Payload and peak & \texttt{mg\_a\_peak\_mem\_gpu5\_r1/summary.json} \\
    Latency breakdown & \texttt{d4\_latency\_m3c\_gpu56\_r1/summary.json} \\
    Streamed-cold smoke & \texttt{p2\_fi\_stream\_cold\_16k\_gpu1\_r1/log\_full.txt} \\
    Weak checks & \texttt{pub\_c\_gemma\_reduced\_gpu01\_r6/summary.json}; \texttt{docs/atlas\_w4\_structure\_rows.jsonl} \\
    \bottomrule
  \end{tabular}
\end{table*}
\endgroup

The complete publication figure suite is generated from tracked artifacts by
\begin{center}
\texttt{uv run python scripts/make\_publication\_figures.py}.
\end{center}
The script emits SVG and PDF files and renders full-size and 3.35-inch PNGs under
\texttt{/tmp/prioritykv-figure-qc} for visual inspection.  The paper compiles
from either the repository root or \texttt{paper/}; the TeX source searches the
appropriate figure directory in both layouts.

\section{Conclusion}
\label{sec:conclusion}

Application-visible structure is a strong retention prior \emph{in a regime we
can now identify in advance}.  On synthetic agent traces it passes 112/120 at a
25\% budget against 1/120 for position-blind eviction and is indistinguishable
from SnapKV-class selection.  On externally authored BFCL multi-turn traces the
same policy scores 0.000 on two different 8B models, while attention selection
remains indistinguishable from FullKV.

The reconciling variable is measurable without running the model: structure can
rank only while protected content is a minority of the context.  Our synthetic
benchmark is 94.9\% filler; real tool-calling prompts are almost entirely
schema.  We therefore recommend the protected fraction as a pre-flight check
rather than structure-aware retention as a default policy, and we give a
budget-relative rule, ADAPT, that degrades to attention selection exactly when
the structural signal runs out.

The surrounding results delimit the claim rather than dilute it: structure
matches but does not exceed FullKV under placement controls, role-aware INT4
placement is not a quality lever, and packed storage is an explicit space--time
trade.  Keeping these hypotheses separate is what makes the retention result
interpretable.  We also publish a correction to one released baseline, so every
column here can be checked against the code.  In operational terms the original intuition survives, narrowed:
when a constrained cache must choose, retain the schema field before the
pleasantry---but first measure whether the context is mostly schema, because then
the rule has nothing left to discard.

\begingroup
\small
\begin{thebibliography}{10}
\setlength{\itemsep}{0.45em}
\setlength{\parsep}{0pt}

\bibitem{vaswani2017attention}
Ashish Vaswani et al.
\newblock Attention is all you need.
\newblock \emph{Advances in Neural Information Processing Systems}, 30, 2017.

\bibitem{kwon2023pagedattention}
Woosuk Kwon et al.
\newblock Efficient memory management for large language model serving with PagedAttention.
\newblock In \emph{Proceedings of the 29th ACM Symposium on Operating Systems
Principles (SOSP)}, 2023.  \href{https://doi.org/10.1145/3600006.3613165}{doi:10.1145/3600006.3613165}.

\bibitem{xiao2024streamingllm}
Guangxuan Xiao et al.
\newblock Efficient streaming language models with attention sinks.
\newblock In \emph{International Conference on Learning Representations (ICLR)},
2024.  \href{https://openreview.net/forum?id=NG7sS51zVF}{OpenReview:NG7sS51zVF}.

\bibitem{zhang2023h2o}
Zhenyu Zhang et al.
\newblock H2O: Heavy-Hitter Oracle for efficient generative inference of large language models.
\newblock \emph{Advances in Neural Information Processing Systems}, 36, 2023.

\bibitem{li2024snapkv}
Yuhong Li et al.
\newblock SnapKV: LLM knows what you are looking for before generation.
\newblock \emph{Advances in Neural Information Processing Systems}, 37, 2024.
\href{https://doi.org/10.52202/079017-0722}{doi:10.52202/079017-0722}.

\bibitem{cai2024pyramidkv}
Zefan Cai et al.
\newblock PyramidKV: Dynamic KV cache compression based on pyramidal information funneling.
\newblock \href{https://arxiv.org/abs/2406.02069}{arXiv:2406.02069}, 2024.

\bibitem{liu2024kivi}
Zirui Liu et al.
\newblock KIVI: A tuning-free asymmetric 2bit quantization for KV cache.
\newblock In \emph{Proceedings of the 41st International Conference on Machine
Learning (ICML)}, PMLR 235:32332--32344, 2024.

\bibitem{ye2025flashinfer}
Zihao Ye et al.
\newblock FlashInfer: Efficient and customizable attention engine for LLM inference serving.
\newblock \emph{Proceedings of Machine Learning and Systems}, 7, 2025.

\bibitem{bai2024longbench}
Yushi Bai et al.
\newblock LongBench: A bilingual, multitask benchmark for long context understanding.
\newblock In \emph{Proceedings of the 62nd Annual Meeting of the Association for
Computational Linguistics}, pages 3119--3137, 2024.
\href{https://doi.org/10.18653/v1/2024.acl-long.172}{doi:10.18653/v1/2024.acl-long.172}.

\bibitem{hsieh2024ruler}
Cheng-Ping Hsieh et al.
\newblock RULER: What's the real context size of your long-context language models?
\newblock In \emph{First Conference on Language Modeling (COLM)}, 2024.
\href{https://openreview.net/forum?id=kIoBbc76Sy}{OpenReview:kIoBbc76Sy}.

\end{thebibliography}
\endgroup

